% Cote 119: offprints, reports, courses, CV and letters (typed and printed part): archive cote 119, site page 170 (PDF page 171).
% typescript. Draft compared: none.
% AI-transcribed by Claude Opus 5.5 (claude-opus-5-5, adaptive thinking, effort high), Message Batches API, batch pilotnext (msgbatch_01A9sdu8SYq23bEtqX9tUncQ), 2026-11-13; not yet checked by eye.
\scanpage{119/170}
\mspage{- 9 -}

Si cela est le cas dans tout l'ouvert  $U$  on dit que  $f$  est \unsure{continûment}
différentiable dans  $U$.

\uline{Définition 4.2.}  (identique à 4.5 de Séries  Entières)

On dit que  $f$  est holomorphe au point  $z_o \in U$  si
\[ \lim_{\substack{u\to 0\\ u\neq 0}} \ (f(Z_o+u) - f(Z_o))/u \ \text{ existe} \]

En notant  $f'(z_o)$  cette limite, on a donc :  $f(Z_o+u)-f(Z_o) = f'(Z_o).u + O(|u|$\unsure{$^\circ$}

En posant  \unsure{$(u =)h + ik$|$u = h + ik$}  on a en comparant avec la relation  (4.1)  le fait
que  $f$  est $\mathbb{R}$-différentiable avec de plus  $a = f'(Z_o)$ et $b = i\, f'(Z_o)$
d'où en remplaçant  $a$ et $b$  par leur valeurs :
\[ \frac{\partial f}{\partial y} = i\,\frac{\partial f}{\partial x} \quad \text{ou encore} \quad \frac{\partial f}{\partial x} + i\,\frac{\partial f}{\partial y} = 0 \ . \]

Réciproquement si  $f$  est  $R$-différentiable et vérifie la relation précédente,
(4.1) montre que  $f$  est holomorphe en  $Z_o$. Si  $f = P + iQ$, les fonctions
$P$ et $Q$  étant réelles on obtient les conditions de Cauchy :
\[ \frac{\partial P}{\partial x} = \frac{\partial Q}{\partial y} \ , \quad \frac{\partial P}{\partial y} = -\frac{\partial Q}{\partial x} \]

\uline{4.3. Théorème de Cauchy}

Si  $f(Z)$  est holomorphe dans un ouvert  $U$  du plan complexe, la
forme différentielle  $f(Z)\, dZ$  est fermée dans  $U$.

\uline{Démonstration :}

D'après la proposition 3.6. il nous suffit de prouver que
$\int_\gamma f(Z)\, dZ = 0$  chaque fois que  $\gamma$  est le bord d'un rectangle contenu (ainsi
que son intérieur) dans  $U$  et ayant ses côtés parallèles aux axes. Soit  $R$
un tel rectangle. Posons a priori  $\int_\gamma f(Z)\, dZ = \alpha(R)$ . Partageons le rectan-%
gle  $R$  en quatre rectangles égaux en partageant chacun de ses côtés en
deux parties égales. Soit  $\gamma_i$  les bords orientés des rectangles obtenus

\mspage{\hand{170}}
